\section{Related Work}
\label{sec:related}

\paragraph{Retrieval units and memory organization.}
Retrieval-augmented generation conditions a language model on an external information store~\cite{rag2020}. The contents and retrieval unit of that store matter as much as the search algorithm. G-Retriever retrieves connected portions of textual graphs for question answering~\cite{gretriever2024}. Its formulation illustrates how retrieval can preserve relations that independent text chunks omit. For manipulation, the corresponding design question is which parts of an interaction should remain connected: a prerequisite, an attempted action, its observed outcome, and a subsequent correction may jointly explain when an experience is useful. RoboRAG therefore treats the graph bank as part of the method. It stores episode graphs and indexes phase-centered subgraphs with evidence references, capability requirements, and versions. Retrieval must select experience that is applicable to the current decision, while retaining the information needed to assess that applicability.

\paragraph{Experience and in-context robot learning.}
RoboICL distinguishes demonstration context from the robot's accumulated interaction memory. Both use observation, action, receipt, and subsequent observation records; anchored memory retains selected earlier interactions~\cite{roboicl2026}. MemER learns to select task-relevant historical keyframes and uses them to produce instructions for a lower-level policy~\cite{memer2026}. These approaches motivate retaining useful history within a limited context budget. Memory Transfer Planning (MTP) is especially relevant to experience reuse across tasks: it stores successful planner code, retrieves examples through instruction similarity, adapts them to the target environment, and uses the adapted examples when replanning after failure~\cite{mtp2025}. RoboRAG examines a different retrieval unit: a phase together with its conditions, feedback, dependencies, and recovery links. Its bank includes observed failed attempts as well as successful continuations, allowing a query to distinguish experiences that share a goal but require different next actions.

\paragraph{Scene graphs, procedural graphs, and task graphs.}
These graph families encode different information. EmbodiedRAG retrieves task-relevant subgraphs from a changing three-dimensional scene representation~\cite{embodiedrag2024}. SafeMem accumulates objects and their relations in a semantic graph and uses this memory to assess action risks~\cite{safemem2026}. Such representations describe the environment. In RoboRAG, a node instead denotes an interaction-phase instance. Object relations appear as conditions or observations attached to that instance; edges describe temporal order, dependencies, conditional transitions, or recovery. This distinction makes explicit whether a record describes what exists, what was attempted, or what should happen next.

Procedural and task graphs are closer precedents. HP-KG constructs a hierarchy of tasks, steps, and actions with an automated multi-agent pipeline. It retrieves procedure nodes and associated subgraphs, then converts them into text for language-model planning~\cite{hpkg2025}. Neural Task Graphs infer a compositional graph from a demonstration and execute it by locating the current node and choosing transitions from visual observations~\cite{ntg2019}. RoboRAG's proposed distinction is the use of recorded phase outcomes and their evidence throughout graph construction, retrieval, adaptation, and execution monitoring. Temporal order alone does not establish a prerequisite or a recovery rule; the annotation procedure must preserve that distinction when converting an episode into reusable experience.

\paragraph{Grounded planning and delegated execution.}
SayCan combines language-model task reasoning with skill affordances~\cite{saycan2022}. Inner Monologue feeds environmental observations and success feedback into subsequent planning~\cite{innermonologue2022}. Code as Policies expresses robot behavior through programs that call perception and control APIs, including feedback loops~\cite{codepolicies2022}. GAVEL provides a particularly close execution comparison: a graph world model validates candidate plans, repairs errors implied by its action semantics, and returns unresolved violations to the language model~\cite{gavel2026}. These methods provide foundations for hierarchical execution and selective language-model invocation. RoboRAG assigns plan construction and substantial revision to a high-level planner, routine phase supervision to a smaller model, and bounded motion execution to local controllers. The question is whether retrieved phase evidence improves decisions at these interfaces, including when to continue, acquire another observation, repair a local plan, or request high-level replanning.

\paragraph{Research position.}
The proposed contribution is a shared representation linking episode annotation, graph-bank retrieval, task-conditioned adaptation, and observed execution. Its evaluation should compare HP-KG-style procedural retrieval, MTP-style reuse and adaptation, graph-based verification such as GAVEL, and linearized versions of the same evidence, with compatible information and controller access. These comparisons test the value of preserving phase, feedback, and dependency relations across the full experience-to-action process.

% Chapter summary: Prior work establishes graph retrieval, procedural reuse, task-graph execution, and selective replanning. RoboRAG studies an evidence-grounded phase graph shared across annotation, retrieval/adaptation, and delegated execution; matched evidence and controller comparisons must establish its value.
